Мышину сказали: <<Эй, Мышин, вставай!>>

Мышин сказал: <<Не встану>>, --- и продолжал лежать на полу.

Тогда к Мышину подошел Кулыгин и сказал:

<<Если ты, Мышин, не встанешь, я тебя заставлю встать>>. <<Нет>>, --- сказал Мышин, продолжая лежать на полу. К Мышину подошла Селезнева и сказала: <<Вы, Мышин, вечно валяетесь на полу в коридоре и мешаете нам ходить взад и вперед>>.

--- Мешал и буду мешать, --- сказал Мышин.

--- Ну знаете, --- сказал Коршунов, но его перебил Кулыгин и сказал:

--- Да чего тут долго разговаривать! Звоните в милицию.

Позвонили в милицию и вызвали милиционера.

Через полчаса пришел милиционер с дворником.

--- Чего у вас тут? --- спросил милиционер.

--- Полюбуйтесь, --- сказал Коршунов, но его перебил Кулыгин и сказал:

--- Вот. Этот гражданин все время лежит тут на полу и мешает нам ходить по коридору. Мы его и так и эдак\dots

Но тут Кулыгина перебила Селезнева и сказала:

--- Мы его просили уйти, а он не уходит.

--- Да, --- сказал Коршунов.

Милиционер подошел к Мышину.

--- Вы, гражданин, зачем тут лежите? --- сказал милиционер.

--- Отдыхаю, --- сказал Мышин.

--- Здесь, гражданин, отдыхать не годится, --- сказал милиционер. --- Вы где, гражданин, живете?

--- Тут, --- сказал Мышин.

--- Где ваша комната? --- спросил милиционер.

--- Он прописан в нашей квартире, а комнаты не имеет, --- сказал Кулыгин.

--- Обождите, гражданин, --- сказал милиционер, --- я сейчас с ним говорю. Гражданин, где вы спите?

--- Тут, --- сказал Мышин.

--- Позвольте, --- сказал Коршунов, но его перебил Кулыгин и сказал:

--- Он даже кровати не имеет и валяется на голом полу.

--- Они давно на него жалуются, --- сказал дворник.

--- Совершенно невозможно ходить по коридору, --- сказала Селезнева. --- Я не могу вечно шагать через мужчину. А он нарочно ноги вытянет, да еще руки вытянет, да еще на спину ляжет и глядит. Я с работы усталая прихожу, мне отдых нужен.

--- Присовокупляю, --- сказал Коршунов, но его перебил Кулыгин и сказал:

--- Он и ночью здесь лежит. Об него в темноте все спотыкаются. Я через него одеяло свое разорвал.

Селезнева сказала:

--- У него вечно из кармана какие-то гвозди вываливаются. Невозможно по коридору босой ходить, того и гляди ногу напорешь.

--- Они давеча хотели его керосином пожечь, --- сказал дворник.

--- Мы его керосином облили, --- сказал Коршунов, но его перебил Кулыгин и сказал:

--- Мы его только для страха облили, а поджечь и не собирались.

--- Да я бы не позволила в своем присутствии живого человека жечь, --- сказала Селезнева.

--- А почему этот гражданин в коридоре лежит? --- спросил вдруг милиционер.

--- Здрасьте-пожалуйста! --- сказал Коршунов, но Кулыгин его перебил и сказал:

--- А потому что у него нет другой жилплощади: вот в этой комнате я живу, в той --- вот они, в этой --- вот он, а уж Мышин в коридоре живет.

--- Это не годится, --- сказал милиционер. --- Надо, чтобы все на своей жилплощади лежали.

--- А у него нет другой жилплощади, как в коридоре, --- сказал Кулыгин.

--- Вот именно, --- сказал Коршунов.

--- Вот он вечно тут и лежит, --- сказала Селезнева.

--- Это не годится, --- сказал милиционер и ушел вместе с дворником.

Коршунов подскочил к Мышину.

--- Что? --- закричал он. --- Как вам это по вкусу пришлось?

--- Подождите, --- сказал Кулыгин и, подойдя к Мышину, сказал: --- Слышал, чего говорил милиционер? Вставай с полу.

--- Не встану, --- сказал Мышин, продолжая лежать на полу.

--- Он теперь нарочно и дальше будет вечно тут лежать, --- сказала Селезнева.

--- Определенно, --- сказал с раздражением Кулыгин.

И Коршунов сказал:

--- Я в этом не сомневаюсь. Parfaitement!